%% introduction --- the book's argument in one sitting: why chaos is worth
%% learning, what it is, and what the five parts do about it.
%%
%% Twin: frontmatter/pl/introduction.tex.

\chapter*{\lblIntroduction}
\phantomsection
\addcontentsline{toc}{chapter}{\lblIntroduction}
\markboth{\lblIntroduction}{\lblIntroduction}

\noindent
The total eclipse of the Sun that crossed Spain on 12 August 2026 was in the
astronomers' tables long before anybody who watched it was born. Its path,
its date and the minute it began had been worked out from the motions of the
Earth and the Moon, and the tables of eclipses run on, calculated the same
way, for thousands of years to come.

Nobody could have told the people waiting under that path, a month
earlier, whether their sky would be clear enough to see it.

Both forecasts are made by the same kind of mathematics. In each case a
computer takes the state of a piece of the world now, applies rules about
how it changes, and steps forward in time. The rules for the air are no
secret: they are the same laws of motion and heat that describe the Moon.
And yet one forecast is good for millennia and the other for about a week.
The difference is not that meteorologists are less careful than
astronomers, and it is not that the air is too complicated to write down.
The difference has a name, a measurement and a precise mathematical
description. It is called chaos, and this book is about it.

\section*{What chaos is}

In everyday speech chaos means a mess. In mathematics it means something
much more interesting, and almost the opposite: a system that follows an
exact rule, with no chance involved anywhere, and still behaves in a way
that cannot be predicted far ahead. Three things have to be true together:

\begin{itemize}[leftmargin=1.4em, itemsep=2pt]
  \item \textbf{the rule is exact} \dash{} the same starting state always
        leads to the same future, with nothing random added;
  \item \textbf{small differences grow} \dash{} two starting states that
        differ by less than anything you could measure end up, after a
        while, completely different;
  \item \textbf{it stays in bounds} \dash{} the motion does not fly off to
        infinity, so the growing differences keep being folded back into
        the same region and never settle down.
\end{itemize}

Each of those is simple. What surprised the scientists who found them
together is how common the combination is. It needs no complicated rule:
the first chaotic system in this book is one line of school algebra, and it
fits on a pocket calculator. It needs no randomness, no crowd and no
mystery. It turns up in dripping taps, swinging pendulums, animal
populations, heart rhythms, the orbits of planets and, most famously, the
weather.

\section*{Why it is worth your time}

You might reasonably ask why anyone who is not a scientist should learn
this. Here are the reasons that persuaded me.

\emph{You live on forecasts.} Weather apps, pension projections, epidemic
curves, economic outlooks, traffic estimates and the confident answers of
computer programs are all predictions made by models. Chaos theory tells you
which of them can in principle be trusted far ahead, which only for a short
time, and what the honest form of a forecast looks like in each case.

\emph{It measures the limit, and the limit is strange.} In a chaotic
system an error in the starting state grows by roughly the same factor every
step. That means improving your measurements a thousandfold does not buy you
a thousand times more forecast; it buys you a fixed extra stretch of time,
and then the fog closes in exactly as before. Once you have seen this worked
out with numbers \dash{} and you will, in Part~II \dash{} you will never
again hear \enquote{we just need better data} in quite the same way.

\emph{It explains why good forecasts come with odds.} If the far future
cannot be computed as a single answer, what can be computed is a spread of
answers and how likely each is. A \enquote{70 per cent chance of rain} is not
a forecaster hedging. It is the correct answer to a question chaos has made
impossible to answer with a yes or no.

\emph{It separates things people usually lump together.} \enquote{We cannot
predict it} can mean that nobody knows the rules, that the rules involve
chance, or that the rules are known exactly and the future is still out of
reach. These are different situations that call for different responses,
and chaos is the third one.

\emph{It is beautiful, and it is universal.} The shapes chaos draws are
fractals, structures with detail at every magnification, and some of its
numbers turn out to be the same in systems that have nothing in common. A
fluid, an electronic circuit and a line of algebra can share a constant that
nobody put there.

\emph{And it is useful.} The same sensitivity that makes chaos
unpredictable makes it steerable with very small nudges, and chaotic
systems can be made to keep in step with each other. Part~V shows both.

\section*{What this book is, and is not}

It is not a textbook, and it will not make you a mathematician. It is also
not a book of pictures with the mathematics taken out. Every idea in it is
the real idea, developed from arithmetic upwards, and every claim a computer
can check has been checked by one; the programs are printed in the book and
you can run them yourself. Where the book states something it does not
prove, it says so and says where the proof is.

It starts from nothing. Part~I (Chapters~\ref{ch:models}
to~\ref{ch:motion}) shows how mathematics describes the world at all: what a
model is, how a rule is stepped forward in time, why feedback stops things
growing for ever, and how motion is described. Part~II
(Chapters~\ref{ch:logistic} to~\ref{ch:bifurcations}) finds chaos in a single
line of algebra, measures how fast small errors grow, turns that into a
single number and a forecasting horizon, and follows the road from order to
chaos. Part~III (Chapters~\ref{ch:lorenz} to~\ref{ch:mandelbrot}) looks at the
shapes chaos makes: the Lorenz attractor, strange attractors, fractals and the
Mandelbrot set. Part~IV (Chapters~\ref{ch:mechanics} to~\ref{ch:computers})
takes it into the world: pendulums and planets, the weather forecast, living
systems and markets, and the computer itself. Part~V
(Chapters~\ref{ch:taming} and~\ref{ch:why}) shows how chaos can be controlled,
and ends with a field guide to judging any prediction you are given.

\section*{The question}

In 1814 the mathematician Pierre-Simon Laplace imagined an intelligence that
knew every force in nature and the exact state of everything at one instant.
For such an intelligence, he wrote, nothing would be uncertain: the future
would be as plain as the past. It is the most confident sentence ever
written about prediction, and for a century it looked like the natural
conclusion of physics.

This book asks whether he was right, and answers with numbers you can
compute yourself. Chapter~\ref{ch:models} begins with a stone and a cup of
tea.
